% Source: https://artofproblemsolving.com/community/c3h1357254_factoring
% Stable identity: aops-topic-1357254-q01
(a) Multiply out; the sum telescopes. For odd $n$,
\[
(x+1)\sum_{j=0}^{n-1}(-1)^jx^{n-1-j}=\sum_{j=0}^{n-1}(-1)^jx^{n-j}+\sum_{j=0}^{n-1}(-1)^jx^{n-1-j}.
\]
The term $(-1)^jx^{n-j}$ for $j\ge1$ cancels the term $(-1)^{j-1}x^{n-j}$ of the second sum. What survives is $x^n$ (from $j=0$ in the first sum) and $(-1)^{n-1}x^0=1$ (from $j=n-1$ in the second sum, because $n-1$ is even). So the product is $x^n+1$. Equivalently, it is the identity $y^n-1=(y-1)(y^{n-1}+\cdots+1)$ with $y=-x$, using $(-x)^n=-x^n$ for odd $n$.

(b) Write $n=km$ with $m$ odd and put $y=x^k$. By (a), $x^n+1=y^m+1$ is divisible by $y+1=x^k+1$. For $x^6+1$ take $k=2$, $m=3$:
\[
x^6+1=(x^2)^3+1=\left(x^2+1\right)\left(x^4-x^2+1\right).
\]
(c) A factor $x-c$ would give a real root $c$, but for even $n$, $c^n+1\ge1>0$. So for even $n$ there is no linear factor; in particular $x^2+1$ cannot be factored over the reals, and $x^{2^k}+1$, which has no odd divisor $m>1$ in its exponent, cannot be split by (b). For example $x^4+1$ does factor, but only as $\left(x^2+\sqrt2x+1\right)\left(x^2-\sqrt2x+1\right)$, with irrational coefficients.
